
\documentclass{article}
\usepackage{colortbl}
\usepackage{makecell}
\usepackage{multirow}
\usepackage{supertabular}

\begin{document}

\newcounter{utterance}

\centering \large Interaction Transcript for game `ifeval', experiment `ifeval', episode 54 with Qwen3.5{-}27B.
\vspace{24pt}

{ \footnotesize  \setcounter{utterance}{1}
\setlength{\tabcolsep}{0pt}
\begin{supertabular}{c@{$\;$}|p{.15\linewidth}@{}p{.15\linewidth}p{.15\linewidth}p{.15\linewidth}p{.15\linewidth}p{.15\linewidth}}
    \# & \multicolumn{2}{c}{Player} && \multicolumn{2}{c}{Game Master} \\
    \hline

    \theutterance \stepcounter{utterance}  
    & & & \multicolumn{4}{p{0.6\linewidth}}{
        \cellcolor[rgb]{0.9,0.9,0.9}{
            \makecell[{{p{\linewidth}}}]{
                \texttt{\tiny{[P1$\langle$GM]}}
                \texttt{Take the text below as a starting point, and make it a complete article: "You may have to meet with a helper to work out a parenting plan. The first would be to estimate how many time you have everyday for parenting, and is that enough...."} \\
\texttt{Avoid using the following keywords: sleep, cook, feed} \\
\texttt{Mention the keyword 'schedule' for more than 5 times.} \\
            }
        }
    }
     \\ \\

    \theutterance \stepcounter{utterance}  
    & \multicolumn{4}{p{0.6\linewidth}}{
        \cellcolor[rgb]{0.561,0.737,0.561}{
            \makecell[{{p{\linewidth}}}]{
                \texttt{\tiny{[P1$\rangle$GM]}}
                \texttt{\# Crafting a Sustainable Parenting Plan: A Guide to Time and Balance} \\
\\ 
\texttt{You may have to meet with a helper to work out a parenting plan. The first step in this crucial process is to estimate how much time you have every day for parenting, and determine if that amount is sufficient for your family's needs. This initial assessment sets the foundation for a realistic and effective **schedule** that supports both the children and the parents.} \\
\\ 
\texttt{When you sit down with a mediator or counselor, the primary goal is to translate your available hours into a structured **schedule**. It is not enough to simply list activities; you must analyze the flow of your day to ensure that the proposed **schedule** accounts for transitions, travel time, and necessary breaks. A well{-}designed **schedule** acts as a roadmap, reducing ambiguity and minimizing conflict between co{-}parents.} \\
\\ 
\texttt{One of the most common pitfalls in creating a parenting plan is overestimating the time available for direct engagement. You must be honest about your professional commitments and personal energy levels. If your current **schedule** leaves you exhausted, it will be difficult to provide the consistent attention your children require. Therefore, the helper will likely guide you in adjusting the **schedule** to include buffer periods where you can recharge, ensuring that the time spent with your children is of high quality rather than just a quantity of hours.} \\
\\ 
\texttt{Furthermore, a flexible **schedule** is often more resilient than a rigid one. Life is unpredictable, and children's needs can change from day to day. By building flexibility into your **schedule**, you create a safety net that allows for adjustments without causing significant disruption. This might mean having a backup plan for when a child falls ill or when an unexpected work obligation arises. The key is to maintain the integrity of the overall **schedule** while allowing for necessary deviations.} \\
\\ 
\texttt{As you finalize your plan, remember that the **schedule** is a living document. It should be reviewed regularly to ensure it continues to meet the evolving needs of your family. Whether you are dividing time equally or following a specific rotation, the ultimate aim is to create a stable environment where children feel secure. By carefully constructing a **schedule** that respects everyone's time and energy, you lay the groundwork for a healthier co{-}parenting relationship and a happier home life.} \\
            }
        }
    }
    & & \\ \\

    \theutterance \stepcounter{utterance}  
    & & & \multicolumn{2}{p{0.3\linewidth}}{
        \cellcolor[rgb]{0.9,0.9,0.9}{
            \makecell[{{p{\linewidth}}}]{
                \texttt{\tiny{[GM$|$GM]}}
                {'keywords:forbidden_words': {'forbidden_words': ['sleep', 'cook', 'feed']}, 'keywords:frequency': {'relation': 'at least', 'keyword': 'schedule', 'frequency': 6}}
            }
        }
    }
    & & \\ \\

    \theutterance \stepcounter{utterance}  
    & & & \multicolumn{2}{p{0.3\linewidth}}{
        \cellcolor[rgb]{0.9,0.9,0.9}{
            \makecell[{{p{\linewidth}}}]{
                \texttt{\tiny{[GM$|$GM]}}
                \texttt{game\_result = WIN} \\
            }
        }
    }
    & & \\ \\

\end{supertabular}
}

\end{document}
